\section{Introduction}\label{sec:intro}

\subsection{The einstein problem in \texorpdfstring{$\R^3$}{R3}}\label{sec:problem}

% ledger: Q1, Q2, Q23, T1, T3, O1, O4, O12 (Figure 1); source 1009.1419 L88; 2509.12216 L249-L253; form 2303.10798 L13, L43
Socolar and Taylor, presenting the three-dimensional version of their hexagonal tile, set out
what remained to be found: ``It is still interesting, however, to search for a single, simply
connected 2D or 3D prototile that forces maximal nonperiodicity by shape alone, or one that does
not permit any weakly nonperiodic tilings'' \cite{ST12}. Writing after the hat, Kaplan restated
the demand from the other side: ``Some people regard this screw motion as uncomfortably close to
translation, and demand a strongly aperiodic 3D monotile, one that admits tilings whose symmetries
never include an infinite cyclic subgroup of any kind. Many of the ideas and algorithms we used to
prove the hat's aperiodicity could be adapted to work in 3D space, but personally I am daunted by
the prospect of deducing the behaviour of any candidate shapes that might be discovered there''
\cite{Kap25}. The solid $\solid$ of \cref{fig:tile} is such a prototile: a closed polyhedral
$3$-ball with rational vertex coordinates that admits tilings of $\R^3$ by congruent copies, and
every tiling it admits has a finite symmetry group (\cref{thm:main,thm:full} below).

\begin{figure}[ht]
  \centering
  \includegraphics[width=\linewidth]{print-preview/figure1-corrected.pdf}
  \caption{The solid $\solid$, with feature bases enlarged $\times 5$ and signed heights
  $\times 80$ for visibility. The seven-cube carrier, feature centres, polarities and relative
  heights are preserved. The outlined panel~$13$ and its inset show the same eight features;
  signed labels give $a$, whose true height is $a/10000$ (positive for a protrusion, negative
  for a recess). The sections illustrate its $+9$ and $-9$ features. The true base side is
  $1/50$; the drawing is not at true feature scale.}
  \label{fig:tile}
\end{figure}

% ledger: Q21; source 2303.10798 L18-L24
The study of aperiodic tilings begins with Wang's tiles and his conjecture that any set of
Wang tiles admitting a tiling admits a periodic one \cite{Wan61}; Berger disproved it with an
aperiodic set of $20{,}426$ Wang tiles and derived the undecidability of the domino problem
\cite{Ber66}. Robinson gave a six-tile aperiodic construction \cite{Rob71} and Penrose a two-tile
one \cite{Pen78}; these counts use different permitted motions and matching conventions from Wang
tiles, for which Jeandel and Rao proved that the minimum is eleven \cite{JR21}. The natural endpoint
of this reduction is a single shape that admits tilings but only non-periodic ones, an aperiodic
monotile or ``einstein'' \cite{SMKGS24}.

% ledger: Q12, Q8, Q21; source 2303.10798 L28-L35; 1903.01158 L11; 2211.15847 L28
Every candidate before the hat conceded something. Gummelt's decagon covers the plane only with
overlaps \cite{Gum96}; Penrose's $1+\varepsilon+\varepsilon^2$ set lets an arbitrarily high
fraction of the area lie in one tile, but the other tiles remain necessary \cite{Pen97}; the
Taylor--Socolar hexagon needs markings that let
non-adjacent tiles exchange information, and a geometric version of the rule needs a disconnected
tile, a tile with cutpoints, or a three-dimensional shape tiling a thickened plane
\cite{Tay10,ST11,ST12}; the Walton--Whittaker tile uses orientational rules \cite{WW23}; and
Mampusti and Whittaker say of their dendritic construction that ``the monotile we present here
is not an einstein in the technical sense'' \cite{MW20}. In a different direction, Greenfeld and
Tao disproved the periodic tiling conjecture with a single tile that tiles $\Z^d$ by translations
only, for a dimension $d$ that ``will be extremely large'' and with a tile that need not be
connected; whether $\Z^3$ admits such a tile they leave open \cite{GT24}, and Greenfeld and
Kolountzakis showed that the tile can be taken connected \cite{GK23}.

% ledger: Q9, Q10, Q25; source 2303.10798 L13, L43; 2305.17743 L45
In the plane the question is settled. The hat, a polykite, is an aperiodic monotile whose tilings
use both the tile and its mirror image \cite{SMKGS24}; the Spectre tiles aperiodically with a
single handedness even when reflections are allowed \cite{SMKGS24b}. Both papers give a
hierarchical, computer-assisted proof, and the hat paper a second, computer-free one. Simplified
proofs and the structure of the tilings followed quickly \cite{AA25,BGS25,BGMM25,Soc23}.

% ledger: Q7, Q6, Q13, Q2, Q24; source 2303.10798 L33, L37; 1009.1419 L14-L16, L30, L69-L72; 1003.4279 L171; 1003.4909 L45-L48
Three dimensions have their own history. The second part of Hilbert's eighteenth problem asked
whether anisohedral polyhedra exist in $\R^3$ \cite{Hil02}; Reinhardt found one \cite{Rei28},
Heesch a planar one \cite{Hee35}. The Schmitt--Conway--Danzer biprism, a convex rhombic biprism
going back to an unpublished 1988 example of Schmitt \cite{Sen96,Rad20}, admits, with reflections
forbidden, only non-periodic tilings: its tilings are periodic layers stacked with a twist, and it
admits a tiling whose symmetry group contains a screw motion; with reflections allowed it admits
a periodic tiling; Socolar and Taylor call the structure ``heterogeneously periodic'' \cite{ST12}. Their own three-dimensional tile enforces the
hexagonal matching rules by shape alone and can be made a topological sphere, but its tilings form
``corrugated slabs of the limit-periodic pattern that stack periodically in the direction
orthogonal to the nonperiodic tiling plane'' \cite{ST11}, which makes it ``a weakly nonperiodic
tiling by Goodman-Strauss's definition'' \cite{ST12}. Fletcher turned \v{C}ul\'ik and Kari's set of
$21$ Wang cubes \cite{CK95} into a single cubic prototile whose $21$ types are $21$ orientations,
but the matching rules become an atlas of legal one-corona patches rather than the shape of the
tile \cite{Fle10}. Kaplan's survey
calls the biprism ``both an existence proof and a cautionary tale about the definition of
aperiodicity in 3D space'' \cite{Kap25}.

% ledger: Q3, Q4, Q5, Q26, T4, T5, N3; source 2303.10798 L33-L34, L64; 2409.15880 L43-L45, L88, L234, L236
That definition is the point on which the three-dimensional problem turns. The hat paper,
``following Mozes'', calls a set of tiles \emph{strongly aperiodic} if it admits tilings but none
with any infinite cyclic symmetry, and \emph{weakly aperiodic} if it admits tilings but none with
a cocompact discrete group of symmetries \cite{SMKGS24,Moz97}. In the Euclidean plane the two
notions coincide for normal tiles \cite[Theorem~3.7.1]{GS16}; in the hyperbolic plane weak
aperiodicity comes readily, as with B\"or\"oczky's monotile \cite{Bor74}, while Goodman-Strauss
constructed a strongly aperiodic set there \cite{GS05}. In $\R^3$ the notions differ: with
reflections forbidden, the biprism is weakly but not strongly aperiodic. We prove the strong form: every tiling by $\solid$ has a symmetry group of
order at most $24$, and a finite group has no element of infinite order (\cref{cor:strong}). Coulbois, Gajardo,
Guillon and Lutfalla use a finer taxonomy in which a tile all of whose tilings have finite
symmetry group is \emph{mildly} aperiodic and ``strongly aperiodic'' is reserved for trivial
stabilizers; in their terms the hat is mildly aperiodic, ``strongly aperiodic monotiles are
unknown in all settings'', and their Question~32 asks for one \cite{CGGL24}. In their terms
$\solid$ is mildly aperiodic and not strongly aperiodic: an explicit tiling by $\solid$ has
symmetry group of order eight, and others have trivial symmetry group
(\cref{sec:aperiodicity,sec:limitperiodic}).

% ledger: O10, O1, O4, N4, T10
$\solid$ is not convex, as the biprism is, and it is not simple: it is a seven-cube chair whose
$24$ exposed unit panels each carry eight square pyramids of height between $1/10000$ and
$12/10000$ (\cref{sec:solid}). It is, however, a closed topological $3$-ball with rational
coordinates, and no matching rule of any kind is imposed: the tilings in \cref{thm:full} are all
the ways congruent copies of the solid can fill space without overlapping.

\subsection{Main result}\label{sec:main}

% ledger: T1-T4, T6, T8, T10, D4, T12; form 2303.10798 L42-L43; 2305.17743 L45
In this paper we prove the following, as a proof submission in the sense explained at the end of
this subsection.

\begin{theorem}\label{thm:main}
The solid $\solid$ of \cref{fig:tile} is a strongly aperiodic monotile of $\R^3$: $\solid$ admits
tilings of $\R^3$ by congruent copies, and no such tiling has a symmetry of infinite order.
\end{theorem}

\begin{theorem}\label{thm:full}
Let $\solid$ be the rational polyhedral solid of \cref{sec:solid}. Then
\begin{enumerate}[label=\textup{(\arabic*)}]
  \item there is a tiling of $\R^3$ by isometric copies of $\solid$;
  \item for every tiling $T$ by isometric copies of $\solid$, $\Per(T)=\{0\}$ and $|\Sym(T)|\le 24$;
  \item every tiling $T$ is homochiral;
  \item in every tiling $T$ each tile lies in a unique infinite hierarchy of $2^n$-scaled chair
    supertiles in registered poses, nested through the eight child poses.
\end{enumerate}
Isometries include reflections. No face-to-face, lattice, common-orientation,
connected-contact-graph or local-finiteness hypothesis is imposed.
\end{theorem}

% ledger: T10, T11, D1, Q3; source 2303.10798 L58
Here $\Per(T)$ is the set of translation vectors $v$ with $T+v=T$ and $\Sym(T)$ the group of
isometries of $\R^3$ that permute the tiles of $T$ (\cref{sec:terminology}). The absence of
hypotheses in \cref{thm:full} is deliberate. A tiling is any family of isometric copies of
$\solid$ with pairwise disjoint interiors whose union is $\R^3$: the copies need not meet
face-to-face, need not sit on a lattice, need not share an orientation, and their contact graph
need not be connected. Local finiteness, which the hat paper notes is automatic for monohedral
plane tilings by closed disks, is proved here from the geometry of $\solid$
(\cref{sec:companions}) rather than assumed.

% ledger: T6, T7, D5, D7, Q7; source 2305.17743 L22; 2303.10798 L52-L54; 1009.1419 L72
Reflections are allowed in \cref{thm:full}, and clause (3) says what they do: nothing. By
longstanding tradition in the tiling literature, ``shapes are considered congruent if they are
equivalent under any Euclidean isometry, including those that reverse orientation''
\cite{SMKGS24}, and the hat's status as a monotile depends on that convention, since every hat
tiling uses both handednesses. $\solid$ needs no such caveat. Its tilings are \emph{homochiral}
in the Spectre paper's sense, ``for every pair $T_i,T_j\in\mathcal{T}$, there is an
orientation-preserving isometry mapping $T_i$ to $T_j$'' \cite{SMKGS24b}: a tiling by $\solid$ is
wholly left-handed or wholly right-handed, and a reflected and an unreflected copy never occur in
the same tiling. By the same paper's definition, a tile that admits only homochiral non-periodic
tilings is a \emph{strictly chiral aperiodic monotile}; $\solid$ is one, in $\R^3$
(\cref{cor:chiral}). The
Schmitt--Conway--Danzer biprism, by contrast, is chiral but admits a periodic tiling once
reflections are allowed \cite{ST12}.

% ledger: T8, T9, N1, N2, A15, R7; source 2305.17743 L43, L45; 2303.10798 L253
Clause (4) is the mechanism. After one ambient isometry every tile of every tiling occupies one
of $24$ proper cubic frames at an integer position (\cref{sec:registration}); the seven-cube
chairs so placed group uniquely into $2$-scaled chairs, those into $4$-scaled chairs, and so on
without end (\cref{sec:hierarchy}). The hierarchy is unique at every level, and that is what
forces non-periodicity, as in the Spectre paper \cite{SMKGS24b}. We do not claim that the
supertiles of one tile exhaust space: as in the hat paper, tilings in which a tile sits at a
corner of its supertile at every level, along ``fault lines'', are not excluded
(\cref{sec:hierarchy}).

% ledger: D4, T12, M1, M2, N8; source README wording (proof submission)
\emph{Proof submission.} The finite parts of the argument, the enumerations of
\cref{sec:computer}, the written geometric lemmas, and the logical assembly from those lemmas to
\cref{thm:full} are theorems checked by the Lean~4 proof assistant, modulo the named compiler
hooks of \cref{app:lean}. The written lemmas concern the tube construction, feature companions,
global coverage and registration, carrier recovery, and small-collar realization. They are new,
written out in full, independently reviewed, and listed by name with the Lean theorems that
prove them in \cref{tab:hypotheses}; there are $19$ of them. Until the development and this
manuscript have been refereed we present the theorem as a proof submission.

\subsection{What an aperiodic monotile is, and how that found \texorpdfstring{$\solid$}{Q}}\label{sec:reading}

% ledger: H1, O5, A15; Lean registration_of_tiling, period_grid_from_registered_poses, period_halving
The proof below follows a reading of the phenomenon that we state first, because it is what
produced the solid. Consider a labelling $x\colon\Z^3\to A$ of the cubic lattice by a finite
alphabet. For a nonzero $p\in\Z^3$, $x$ is $p$-periodic exactly when it factors through
$\Z^3\to\Z^3/\langle p\rangle$; this quotient is infinite. For tilings by $\solid$ the connection
to labellings is
registration (\cref{sec:registration}): after one ambient isometry every tile of every tiling
sits in one of $24$ proper cubic frames at an integer position, and, because $\solid$ has no
self-isometry, every translational period of the tiling is then an integer vector
(\cref{sec:aperiodicity}). We use only this direction; the paper does not identify the periods of
a tiling with those of a labelling in any packaged form. For $\solid$, we prove existence and
exclude every nonzero translational period directly, using registration and period halving. The
obstruction is productive: the local pieces do fit together globally; what fails is fitting them
together periodically.

% ledger: H3, Q11; source 1608.07165 L80, L308; 2305.17743 L43
A hierarchy certifies aperiodicity only when it is forced. Grouping the tiles of a tiling into
supertiles proves nothing by itself: on the cubic grid an observer may group cells into
$2\times2\times2$ blocks in eight ways, and the grouping adds information the tiling does not
carry. What proves non-periodicity is a grouping that the tiling determines and that every
symmetry of the tiling therefore preserves. Then a symmetry of the tiling is a symmetry of the
grouped tiling at every level, so a translation that was a symmetry would survive coarsening to
every scale; if each coarsening halves it, it was never there. This is the recognizability,
unique-composition and period-halving mechanism used here \cite{GS18,SMKGS24b}, and we use it as
a design target.

% ledger: H4; source packets P3-P5 SOURCE_AUDIT (structural interpretation); history/cascade/findings.md L44
With marked tiles the rule that decides which neighbours are legal is external to the pieces; a
geometric monotile has to put that rule into what can physically fit. For a marked system with a
forced hierarchy to become one bare solid, three things must hold: a solid realizing the marked
system exists; every tiling by the solid decodes into a legal configuration of the marked
system; and a period of the tiling is a period of the decoded configuration. The second is the
one that is easy to lose. It is not enough that the intended tilings survive: every tiling
admitted by the bare solid must decode into the intended admissible marked system. One shape does
not mean one role: copies of a
single solid occupy different positions in the forced structure, and the diversity that a set of
marked tiles carries in its catalogue is carried here by the relations between congruent copies.

% ledger: H5, R7; Lean parent_atlas_eq_fine
This reading yields a finite test. Write the legal contacts between neighbouring copies of the
solid as a language of contacts, and decode a tiling into the tiling by its parents. The halving
argument needs the parent tiling to be again a legal tiling of the same solid, that is, the
language of contacts among decoded parents must lie inside the language of contacts among the
tiles themselves. Equality of the two languages is the form we check, and it is sufficient for
the argument, not necessary: a strictly smaller decoded language would also do. When the decoded
language is strictly larger the argument fails, and whether a periodic parent tiling then exists
must be checked case by case; for the first bare solid we built, it did (\cref{sec:found}). For
$\solid$ the test is one finite identity, verified by machine: the halved parent atlas equals the
fine atlas of $44$ legal contacts (\cref{sec:designtest}).

% ledger: H2, H6, H11 (pointer sentence only; framework vocabulary confined to sec:found)
The reading was reached with the Six Birds emergence calculus \cite{Tsi26F6,Tsi26F4,Tsi26NDO};
\cref{sec:found} quotes its statement of the phenomenon and relates it, clause by clause and at
each clause's level of verification, to the theorems proved here. Nothing in that section is used
in the proofs.

% ledger: H11, N7, O8, R10, A7, A11, A15, R6, R7, T13; source 2303.10798 L51; figure roadmap
The proof is the discharge of that reading, and \cref{fig:roadmap} shows its shape. Existence
comes from a substitution: eight rotated copies of the chair form a doubled chair, and iterating
the dissection gives a registered tiling of the chair which the features accept
(\cref{sec:finding}). Universality comes from the features: each pyramid forces a companion
pyramid of the opposite polarity on a neighbouring tile, the companion forces the neighbour's pose
into a discrete set, and an exhaustive census reduces that set to the $44$ legal contacts
(\cref{sec:companions}); from there every tiling is registered (\cref{sec:registration}), every
registered tiling has a unique parent tiling at every level (\cref{sec:hierarchy}), and a
translational period would have to halve at every level, hence vanish (\cref{sec:aperiodicity}).
As in the hat paper, the computer-assisted census is needed only to show that all tilings follow
the hierarchy; it is not needed to show that a tiling exists.

\begin{figure}[ht]
  \centering
  \resizebox{\linewidth}{!}{% fig_roadmap.tex — Fig. 1.2, proof roadmap (FIGURES.md 1.2; hat Figs 1.2/1.3 role).
% Two flowcharts: existence (top) and universality (bottom). Each box carries its section and
% the CLAIMS_LEDGER.md ids it stands on; tiers are those of the ledger rows at the pinned commit.
\begin{tikzpicture}[
  node distance=4mm and 5mm,
  box/.style={draw, rounded corners=1.5pt, align=center, font=\footnotesize, inner sep=3pt,
              minimum height=9mm, text width=27mm},
  goal/.style={box, thick, fill=black!6},
  lab/.style={font=\scriptsize\itshape, text=black!60},
  arr/.style={-{Latex[length=2mm]}, thick}
]
  % --- existence ---
  \node[lab] (exlab) {Existence (\cref{sec:finding}): a tiling of $\R^3$ by $\solid$};
  \node[box, below=2mm of exlab.west, anchor=north west] (e1)
    {\S\ref{sec:substitution} eight child poses partition $2\carrier$ \\ (O8)};
  \node[box, right=of e1] (e2)
    {\S\ref{sec:substitution} the $44$-contact atlas, $24$ frames \\ (R2, R3)};
  \node[box, right=of e2] (e3)
    {\S\ref{sec:existence} registered baseline tiling \\ (R10)};
  \node[box, right=of e3] (e4)
    {\S\ref{sec:existence} small-collar realization \\ (R10, A16)};
  \node[goal, right=of e4] (e5)
    {\cref{thm:full}(1) \\ a tiling exists (Lean spine; \cref{app:lean})};
  \draw[arr] (e1) -- (e2); \draw[arr] (e2) -- (e3); \draw[arr] (e3) -- (e4); \draw[arr] (e4) -- (e5);

  % --- universality ---
  \node[lab, below=9mm of e1.south west, anchor=north west] (unlab)
    {Universality (\cref{sec:companions,sec:registration,sec:hierarchy,sec:aperiodicity}): every tiling by $\solid$};
  \node[box, below=2mm of unlab.west, anchor=north west] (u1)
    {\S\ref{sec:companions} features force companions \\ (A1--A9)};
  \node[box, right=of u1] (u2)
    {\S\ref{sec:registration} only registered mates; the census \\ (A10--A12)};
  \node[box, right=of u2] (u3)
    {\S\ref{sec:registration} registration, \cref{thm:registration} \\ (A13--A15)};
  \node[box, right=of u3] (u4)
    {\S\ref{sec:hierarchy} unique parent, \cref{thm:parent} \\ (R4--R6)};
  \node[box, right=of u4] (u5)
    {\S\ref{sec:hierarchy} coarsening; the hierarchy \\ (R7, T8)};
  \draw[arr] (u1) -- (u2); \draw[arr] (u2) -- (u3); \draw[arr] (u3) -- (u4); \draw[arr] (u4) -- (u5);

  \node[goal, below=of u5] (u6)
    {\cref{thm:full}(2): period halving, $\Per(T)=\{0\}$, $|\Sym(T)|\le 24$ \\ (T13, T2, T3)};
  \node[goal, left=of u6] (u7)
    {\cref{thm:full}(3): homochiral \\ (T6)};
  \node[goal, left=of u7] (u8)
    {\cref{thm:full}(4): unique hierarchy \\ (T8)};
  \draw[arr] (u5) -- (u6);
  \draw[arr] (u3.south) |- ([yshift=2mm]u7.north) -- (u7.north);
  \draw[arr] (u5.south west) -- (u8.north east);
\end{tikzpicture}
}
  \caption{Roadmap of the proof. Each box names its section and the rows of the claims ledger
  (\cref{app:lean}) it rests on.}
  \label{fig:roadmap}
\end{figure}

\subsection{Computer assistance and verification}\label{sec:computer}

% ledger: M1-M5, M7, T12, Q14, Q20, R2, R4, R6, R7, A11, L1; source 2303.10798 L51, L158; 1808.07768 L51, L253; 1506.06492 L94
Several steps of the proof are exhaustive finite enumerations, listed in \cref{tab:census}: the
$2{,}388$ candidate contact poses of which exactly $44$ are legal; the $33$ first shells and $28$
parent conflicts behind the unique-parent theorem; the $697\to116\to44$ parent contacts of the
coarsening; the companion census of $6{,}862$ mate poses reduced to $44$ through $299{,}975$
collision boxes; and the $168$-label substitution with its primitivity exponent and modular
coincidence depth. The companion census, the largest of these, was checked by two independent
Python implementations. Every enumeration is stated as a
theorem in Lean~4 and checked by its kernel: those that decide directly by \lean{decide}, and
those too large for the kernel's evaluator by \lean{native\_decide}, which adds, per theorem, one
named compiler trust hook that we list in \cref{app:lean}. No floating-point number enters any
checked statement: certificates are integers and exact rationals, the Python replay uses exact
arithmetic, and Lean decides integer literals. The logical assembly from the enumerations and the
named written lemmas to \cref{thm:full} is itself a Lean theorem, and the written lemmas of
\cref{sec:solid,sec:finding,sec:companions,sec:registration} are themselves Lean theorems, so that
the main theorem is kernel-checked modulo the named compiler hooks listed in \cref{app:lean}; the
written proofs remain as exposition. Negative
controls (must-fail variants of the inputs and of the formal statements, each checked to fail as
expected) accompany the positive build. A one-command replay of the enumerations from the Python
standard library runs in about a minute (\cref{app:data}). The only comparable formalization we
know of is Myers's staging repository for the hat and Spectre, still in progress
\cite{Mye24lean,Lean4}. \emph{Code.} The repository is given in \cref{app:data}; its release tag is
\texttt{paper-v1} (intended; tag creation and licensing remain pending owner actions).

\subsection{Terminology}\label{sec:terminology}

% ledger: D1, D2, D3, D9; source 2303.10798 L25, L57-L58, L61; 1003.4909 L25
Terminology follows Grünbaum and Shephard as the hat paper uses it \cite{GS16,SMKGS24}, with
``closed topological disk'' replaced by ``closed topological $3$-ball''. A tile is a closed set;
a tiling by a set of tiles is a collection of images of tiles under isometries with pairwise
disjoint interiors whose union is the whole space; the set, or the single tile, \emph{admits} the
tiling. A tiling is monohedral if all its tiles are congruent, congruences including mirror
reflections, and locally finite if every ball meets only finitely many tiles. A \emph{monotile}
is a tile admitting a monohedral tiling; adapting the hat paper's reservation of the terms to the
present setting, an \emph{aperiodic monotile} or \emph{einstein} here is a closed topological
$3$-ball admitting tilings but only non-periodic ones, by geometry alone, without any
non-geometric matching rule constraining tile adjacencies. \emph{Registered} tilings, the \emph{atlas} of legal contacts,
and the \emph{features}, \emph{panels} and \emph{carrier} of $\solid$ are our own terms, defined
in \cref{sec:solid,sec:finding}; a \emph{supertile of level $n$} is a $2^n$-scaled chair in a
registered pose (\cref{sec:hierarchy}). Here a patch means a finite subfamily of a packing;
unlike the hat paper's disk-patch convention \cite{SMKGS24}, we impose no topological-ball
condition. A $1$-corona consists of the tiles touching the chosen tile \cite{Fle10}. The notation
is collected in \cref{app:notation}, \cref{tab:notation}.

% ledger: D1 (symmetry), Q3, T5, Q5, O5; source 2303.10798 L34, L64; 2409.15880 L43-L45, L73-L76
We quote the hat paper's definitions. ``The symmetry group of a tiling is the group of those
isometries that act as a permutation on the tiles of the tiling. A tiling is weakly periodic if
its symmetry group has an element of infinite order \dots\ A tiling is strongly periodic if the
symmetry group has a discrete subgroup with cocompact action on the space tiled \dots\ A set of
tiles (or a single tile) is weakly aperiodic if it admits a tiling but does not admit a strongly
periodic tiling, and strongly aperiodic if it admits a tiling but does not admit a weakly periodic
tiling.'' And: ``Following Mozes [Moz97], we say a set of tiles is strongly aperiodic if it admits
tilings but none with any infinite cyclic symmetry'' \cite{SMKGS24,Moz97}. Since a finite group
has no element of infinite order, ``strongly aperiodic'' in this sense is what
\cref{thm:full}(2) gives. In the finer taxonomy of Coulbois, Gajardo, Guillon
and Lutfalla \cite[\S2.1]{CGGL24} a tile all of whose tilings have finite symmetry group is
\emph{mildly} aperiodic, and \emph{strongly} aperiodic is reserved for trivial stabilizers; in
their terms $\solid$ is mildly aperiodic, and it is not strongly aperiodic in that stricter
sense, since an explicit tiling has symmetry group of order eight (\cref{sec:aperiodicity}).
They also distinguish the symmetry group of a tiling from the
stabilizer of its family of placements; the two coincide when the tile has no self-isometry, as
$\solid$ has none (\cref{sec:solid}), and we write $\Sym(T)$ throughout.

% ledger: D5, Q10; source 2305.17743 L22, L43
From the Spectre paper we take, verbatim, that ``a monohedral tiling $\mathcal{T}$ is a
homochiral tiling if for every pair $T_i,T_j\in\mathcal{T}$, there is an orientation-preserving
isometry mapping $T_i$ to $T_j$'' and that a strictly chiral aperiodic monotile is ``a tile that
admits only homochiral non-periodic tilings''; and that ``a set of tiles is hierarchical if, in
every tiling admitted by those tiles, every tile is nested within an infinite hierarchy of
ever-larger supertiles. If these hierarchies are uniquely determined, then the tilings that
contain them must be non-periodic'' \cite{SMKGS24b}. Fault lines are defined in
\cref{sec:hierarchy}, following the hat paper.
